\MajorSection{الملاحظة الأولى}

\MinorSection{الحاج عبدو الإنسان}

عرفت الحاج عبدو منذ سنوات أكثر مما أحب أن أسجّل. وهو، كما يُعتقد، من دارابگرد في ولاية يزد؛ وكان يؤثر دائمًا أن يسمّي نفسه «الهيچمكاني»، وهو لقب مازح يعني «من لا دار، من لا مكان». لقد سافر بعيدًا وفي جهات شتى بعينين مفتوحتين، كما يظهر من «أبياته المزدوجة». وإلى سليقته الطبيعية وموهبته في تعلّم اللغات، أضاف حصيلة من قراءات متفرقة متنوعة: شذرات من الصينية والمصرية القديمة، والعبرية والسريانية، والسنسكريتية والبراكريتية، واللغات السلافية، ولا سيما الليتوانية، واللاتينية واليونانية، بما فيها الرومية الحديثة، والأمازيغية واللهجة النوبية والزندية والأكدية؛ فضلًا عن الفارسية، لسانه الأم، والعربية، لغة الدرس الكلاسيكية. ولم يكن يجهل تلك العلوم ذات الأسماء المنتهية بـ«ـولوجيا»، ولا انتصارات الكشف العلمي الحديث. وخلاصة القول أن ذاكرته كانت غنية المخزون؛ وكان يملك كل موهبة إلا موهبة استعمال مواهبه.

لكن لم يخطر لأحد أنه «يخطب ودّ ربة الشعر»، إن تكلمنا بأسلوب القرن الماضي. حتى أقرب الناس إليه كانوا يجهلون أن له سرًا يخفيه في خزانته: قصيدته، أو أبياته المزدوجة. وقد أسرّ إليّ بهذا السر حين التقينا آخر مرة في غربي الهند؛ وأنا أتعمد الإبهام في تحديد الموضع. وهو يفعل ذلك، أمسك بيده مجد ذقنه الطويل الأشيب، موجّهًا أطرافه نحوي، كأنه يقول مع ملك الجزيرة:

\begin{SourceQuotation}
في لحيتي مسّة من الشتاء؛\\
علامة على أن الآلهة ستحميني من الطيش.
\end{SourceQuotation}

ومع ذلك، بدت عينه النافذة، الصافية كالعقيق، معترضة على حجة العمر. وكانت المخطوطة مكتوبة بأسوأ خط «شكسته»، أو الخط الجاري؛ وحين أخذتها معي أبى الكاتب أن يتجشم عناء نسخ تلك الكتابة الرديئة.

كنا نحن أصدقاءه القدامى نخاطب الحاج عبدو منذ زمن بلقب «نبينا»؛ وسيرى القارئ أن للحاج رسالة يريد تبليغها، أو أنه يعتقد أن له رسالة. ومن الواضح أنه يطمح إلى الدعوة إلى إيمان خاص به: صيغة شرقية من المذهب الإنساني، ممتزجة بالعادة الذهنية الشكّية، أو العلمية كما نقول الآن. وأحسب أن الفيلسوف قد يقبل هذا الدين الذي لا تعدو عبادة الأوثان والهندوسية والوثنية، واليهودية والمسيحية والإسلام، أن تكون أجزاء منه؛ فهو يعبد بتفانٍ لا ينقسم قضية الحقيقة المقدسة، والحقيقة لذاتها، لا للخيرات التي قد تجلبها. وهذا الاعتقاد مقبول على السواء عند الجهل الصادق وعند أعلى مراتب دراسة الطبيعة.

وينمّي الحاج، مع كونفوشيوس، ما سمّاه شتراوس «الحس السليم الصارم للبشر»؛ فيما تكون سيادة النظام فقرة من «ناموسه الأعلى». ويتتبّع من أبسط البدايات ذلك الإدراك البشري الذي يُدعى «الإيمان»، وقد بلغ من العموم ما يقارب الإطلاق: ذلك الحس بالألوهية، \textenglish{sensus Numinis}، الذي أصبح الآن، بالوراثة أو الانتقال، عامًا إلا عند من يحملون أنفسهم على معارضته. ومن الواضح أنه يرى في هذا الاتفاق العام للبشر قدرًا من الإلهية؛ إذ اكتشف لنفسه في الأصل ألوهية، إن لم يكن قد خلقها. ولا يصيح مع مسيح نوفاليس: «يا أطفال، ليس لكم أب»؛ وربما شارك رينان قوله: «عالم بلا إله أمر مروّع!»

لكنه يدرك عدم التوافق بين اللامتناهي والمحدّد؛ بين كائن يحب ويفكر ويكره، وبين «فعل محض»، \textenglish{Actus purus}، يُقال إنه غيور وغضوب ومنتقم، وبين «أزلي يعمل في سبيل البر». وأمام التناقضات التي لا تنتهي، الناشئة من فكرة إله شخصي، ومن تركيب العناية الإلهية ومفهومها، يلجأ صاحبنا اللاأدري إلى الشعور بمجهول وبما لا يمكن معرفته. ويعترض على تنوع الصور التي لا تُحصى، والتي يتخذها إدراك «العلة المعلِّلة»، \textenglish{Causa Causans}، وهو اسم غير صائب؛ كما يعترض على ذلك التبنّي العقلي لقضايا عامة قابلة للصياغة الواضحة، غير قابلة للإثبات، والذي نسمّيه الاعتقاد.

وينظر بعين محايدة إلى التنوع الذي لا ينتهي من المذاهب التي يتمسك بها، بثقة واكتفاء متساويين، رجال متساوون في القدرة والصدق. لقد سئم الطواف بالعالم، والعثور على كل شعب صغير متزوج بآرائه الخاصة، يدّعي احتكار الحقيقة، ويعد الآخرين جميعًا مخطئين، ويثير نزاعات تتناسب شدتها وحدتها وسمّيتها عكسيًا مع أهمية الأمر المتنازع عليه. وقد علّمته ملاحظة بالغة النشاط والنفاذ أن كثيرًا من هذه العائلات المتخاصمة، ولا سيما ما كان منها من دم واحد، متكافئة في عمليات الإدراك والتأمل العقليتين؛ وأنها، في شؤون العالم المرئي العامل، لا يقرّ أحد بتفوّق بعضها على بعض. أما في المسائل الغامضة المتعلقة بالإيمان المحض، والتي لا تقبل دليلًا مباشرًا حسيًا، فإن واحدًا من كل مئة سيدّعي الصواب، ويتهم التسعة والتسعين الآخرين بالخطأ من غير تواضع.

ومن ثم يسعى إلى اكتشاف مذهب يثبت أنهم جميعًا على صواب، وجميعًا على خطأ؛ ويوفّق بين اختلافاتهم، ويوحّد عقائد الماضي، ويفسر الحاضر، ويستبق المستقبل بتطور متصل لا ينقطع. ويكون ذلك أيضًا بعملية لا تقوم على النفي والتمييز، بل على قدر بالغ من الإثبات والبناء. ولست مدعوًا إلى الجلوس في مقعد الحكم، لكن يسعني أن أقول إن نجاح المحاولة سيكون أمرًا عجيبًا. ومذهب كهذا سيكون شاملًا لكل شيء، إذ لا يحده مكان أو زمان أو شعب؛ وسيكون مبدؤه ممتدًا امتداد المادة نفسها، ومن ثم أبديًا. وفي أثناء ذلك يرضي نفسه؛ وهذا هو الأمر الأساسي.

عرّف دارسو الميتافيزيقا في السنوات الأخيرة إساءة استعمال علمهم بأنها «مورفولوجيا الرأي الشائع». ويقولون إن الباحثين المعاصرين انشغلوا أكثر مما ينبغي بالتأمل الباطني؛ وصارت أعمالهم مجرد دراسات تجمع وظائف الأعضاء إلى السيرة الفردية، وأهملوا إهمالًا كبيرًا دراسة المتوسطات. إذ يقول لاروشفوكو: «معرفة الإنسان بوجه عام أيسر من معرفة إنسان بعينه»؛ وفي موضوع بهذا الاتساع لا بد أن تكون كل وجهة نظر أحادية الجانب.

لكن هذا ليس نهج الشرقيين. فما زال عليهم أن يعالجوا المسائل الكبرى بالقياس إلى الكون، \textenglish{ex analogiâ universi}، بدل القياس إلى الإنسان، \textenglish{ex analogiâ hominis}. وعليهم أن يتعلموا أساس علم الاجتماع: ذلك الاقتناع الفلسفي بأن البشر ينبغي أن يُدرسوا كلًا عضويًا، لا كومة من أفراد. ولهذا تحتقر روح العصر، \textenglish{Zeitgeist}، أو التطور التاريخي للوعي الجمعي للعصر، الرأي البالي القائل إن المجتمع، أي الدولة، مقيّد بالواجبات الأخلاقية نفسها التي تقيّد المواطن البسيط. ولهذا أيضًا ترى أن «روح الإنسان، لكونها من جوهر متساوٍ ومتجانس، تفترض عادة وتتخيل في الطبيعة مساواة وتجانسًا أكثر مما فيها في الحقيقة».

لقد خضعت المسيحية والإسلام للاختبار خلال القرون الثمانية عشر والاثني عشر الماضية. وكانتا شديدتي الحماسة لنشر الدعوة، ومع ذلك لا تضمان إلا عُشر الجنس البشري ونصف عُشره. ويُرجع الحاج عبدو التقدم البطيء غير المرضي لما يسميه أتباعهما «الحقائق الخالصة» إلى ما فيها من نقائص أصيلة. فكلتاهما تعرض ثوابًا على مجرد الاعتقاد، وعقابًا على مجرد عدم الاعتقاد؛ والثواب والعقاب، بالمناسبة، غير متناسبين تناسبًا بالغًا. وهكذا تردّان كل شيء إلى مقياس أنانية غير مهذبة كثيرًا؛ وتصبح آثارهما المفسدة للأخلاق أوضح مع كل عصر يتقدم.

لا يطلب الحاج عبدو إلا الحقيقة؛ الحقيقة بالقدر الذي يستطيع الإنسان فهمه في الطور الراهن من تطوره. وهو يأنف أن يقرن المنفعة، كما فعل بيكون، في الأورغانون الجديد، الكتاب الأول، الحكمة ١٢٤، بذلك «النور الجاف النقي للمفاهيم الحقيقية»، \textenglish{lumen siccum ac purum notionum verarum}؛ وبيكون هو الكاهن الأكبر للعقيدة الإنجليزية، «الحس العملي الغليظ»، \textenglish{le gros bon sens}. ويبدو أنه يرى الضرر الذي ألحقه بمجموع الفكر ذلك التطيّر بالمعرفة البعدية، وعبادة «الوقائع»، وتأليه التركيب. وأخيرًا جاءت الطريقة الطائشة التي «حرّر» بها لوك «الفلسفة من كابوس الأفكار الفطرية». فمثل لوثر وقادة الثورة الفرنسية الكبرى، قطع الصلة بالماضي، وألقى من السفينة كامل حمولة التراث الإنساني. وكانت النتيجة حركة هائلة للعقل نحب أن نسميها تقدمًا، مع أنها كثيرًا ما كانت نكوصًا؛ مع نمو عظيم للأنانية، ناشئ عن الشعور المدلل بالشخصية الفردية.

يأسف الحاج للأهمية المفرطة المعطاة لحالة مستقبلية محتملة؛ ويرى فيها منبهًا نفسيًا وحلم يقظة، يخلّ انتكاسه ورد فعله بحياة اليقظة. وقد تبدو حالته متواضعة نثرية لمن ارتفعوا بأبخرة الخيال، وباحتساء مسكر روحي لا يختلف عن المسكر الجسدي في أنه طلب لسعادة مثالية. لكنه أحكم من أن يثبت وجود عالم آخر أو ينفيه. وليس ثمة اتفاق للبشر على حياة وراء القبر، ولا رأي جامع اعتُقد «دائمًا، وفي كل مكان، ومن الجميع»، \textenglish{semper, et ubique, et ab omnibus}. والقوى العقلية، الإدراك والتأمل، صامتة في هذه المسألة؛ لا تشهد بوقائع، ولا تقدّم برهانًا. حتى الحس الغريزي لجنسنا أبكم هنا. قد نؤمن بما نتعلمه؛ لكننا لا نستطيع أن نعرف شيئًا. ولذلك يريد أن ينمّي تلك الحال المتقبلة التي، إذ تسير في ظل أحداث عظيمة، تقود إلى أسمى الغايات: تطور الإنسانية. وتعليق الحكم عنده مذهب.

لقد أنجز الإنسان كثيرًا خلال القرون الثمانية والستين التي تمثل تاريخه. ويقوم هذا على افتراض أن الإمبراطورية المصرية الأولى، التي أعقبت ما قبل التاريخ، بدأت سنة ٥٠٠٠ قبل الميلاد، وانتهت سنة ٣٢٤٩ قبل الميلاد. وكانت القديمة، في مقابل الوسطى والجديدة والمتأخرة؛ وضمت الأسر من الأولى إلى العاشرة؛ وكانت عصر الأهرام، البسيطة والمتينة والعظيمة في آن واحد. وحين يؤكد مادح الماضي أن الحضارة الحديثة لم تتحسن في شيء عمّا عند هوميروس وهيرودوت، فإنه ينسى أن كل تلميذ مدرسة معجزة من التعلم إذا قورن بإنسان الكهف وبجنس العصر الحجري القديم. وكما كان الماضي سيكون المستقبل.

نظرة الحاج إلى الحياة هي نظرة الصوفي، مع المسحة المعتادة من التشاؤم البوذي. وقد انتقل حزن الوجود العميق، الذي كثيرًا ما غنّاه الشاعر الشرقي الحالم، إلى العقل الأوروبي العملي. حتى الفرنسي الخفيف يهمس:

\begin{SourceQuotation}
أنا، أنا، أحني رأسي كل يوم أكثر؛\\
وأمضي، وقد بردت تحت هذه الشمس المرحة؛\\
وسأرحل قريبًا، في قلب الاحتفال؛\\
من غير أن ينقص هذا العالم الفسيح المضيء شيء.
\end{SourceQuotation}

لكن حاجنا ليس عدميًا بمعنى «لا شيء البتة» في قصيدة هود؛ أو، كما يعبّر الأمريكي: «لا شيء جديد، ولا شيء صحيح، ولا أهمية للأمر». إن نواحه سليم على قِصر الحياة ومصائبها؛ لأنه يجد المخلوقات كلها:

\begin{SourceQuotation}
تقيس العالم بـ«أنا» هائلة.
\end{SourceQuotation}

ويذكّرنا بالقديس أوغسطين، التأملات، الفصل ٢١: «هذه الحياة، الحياة البائسة، الزائلة، غير اليقينية، الشاقة، الدنسة؛ سيدة الشرور، وملكة المتكبرين، المليئة بالمصائب والأخطاء... تنفخها الأخلاط، ويملؤها الطعام، ويُهزلها الصيام، وتُرخِيها الملاهي، وتستهلكها الأحزان؛ يضيّقها القلق، ويبلّدها الأمان، وتنفخها الثروة وتقلّبها. يخفضها الفقر، ويرفعها الشباب، وتحنيها الشيخوخة، ويكسرها الإلحاح، ويذلّها الحزن. وبعد هذه الشرور كلها يأتي الموت المستعر». لكن الحاج ربما قرأ «المبارك»، \textenglish{benedicta}، بدل «المستعر»، \textenglish{furibunda}.

ويجد الحاج عبدو، مع الكاردينال نيومان، أحد أمجاد عصرنا، أن «نور العالم ليس إلا لفافة النبي، مليئة بالنواح والحداد والويل». ولا أستطيع الامتناع عن اقتباس هذا المقطع الجميل كله، ولو من أجل استنتاجه الأعرج الضحل وحده: «أن تتأمل العالم طولًا وعرضًا، وتاريخه المتنوع وأجناس البشر الكثيرة، وبداياتهم وأقدارهم وتباعدهم بعضهم عن بعض ونزاعاتهم؛ ثم طرائقهم وعاداتهم وحكوماتهم وصور عبادتهم؛ ومشاريعهم ومساراتهم التي لا غاية لها، وإنجازاتهم ومكتسباتهم العارضة؛ والختام العاجز لوقائع طال أمدها؛ والعلامات الواهنة المتقطعة على تدبير مشرف؛ والتطور الأعمى (!) لما ينتهي إلى قوى أو حقائق عظيمة؛ وسير الأشياء كما لو أنها تنطلق من عناصر بلا عقل، لا نحو علل غائية؛ وعظمة الإنسان وضآلته، وغاياته البعيدة المدى وقصر مدته؛ والستار المسدل على مستقبله، وخيبات الحياة، وهزيمة الخير، ونجاح الشر، والألم الجسدي والعذاب النفسي، وشيوع الخطيئة وشدتها، وعبادات الأوثان المتفشية، والفساد، واللادينية الكئيبة اليائسة؛ وتلك الحال للجنس كله التي وصفها قول الرسول بذلك الهول وتلك الدقة: بلا رجاء وبلا إله في العالم؛ كل هذا مشهد يدوّخ ويفزع، ويفرض على العقل إحساسًا بسر عميق لا حل إنسانيًا له على الإطلاق». ولهذا يفترض ذلك الكاتب الجدير بالإعجاب «كارثة أصلية مروّعة»؛ فتغدو عنده العقيدة البغيضة التي تسمّى لاهوتيًا «الخطيئة الأصلية» يقينية تقريبًا بقدر يقين «وجود العالم ووجود الله». وكذلك تصر «قائمة العقائد» الخاصة بأكثر الكنائس المسيحية تحررًا على فساد الإنسان، وعلى «الضرورة المطلقة لفعل الروح القدس في تجديد الإنسان وتقديسه».

ولكن ماذا لدينا هنا؟ إن «الكارثة الأصلية» إما أن الله أحدثها، وإما أنها نشأت بغير إذن الله؛ وفي الحالتين يُهبط بالله إلى الإنسان. إنها المعضلة القديمة التي يتمثل طرفاها في صفتي الخير والعلم المحيط غير القابلتين للتوفيق في الخالق المفترض للخطيئة والألم. فإن أمكن حمل إحداهما عليه، لم يمكن حمل الأخرى على الموضوع نفسه. وخير وأحكم بكثير ذلك التفسير الشعري لصاحب المقالة، والذي يبدو أنه صار محتقرًا لأنه كان مذهبًا رائجًا في أيام الشاعر الحكيم:

\begin{SourceQuotation}
الطبيعة كلها ليست إلا صنعة...\\
وكل تنافر انسجام لم يُفهم؛\\
وكل شر جزئي خير كلي.\\
المقالة، الأبيات ٢٨٩ إلى ٢٩٢.
\end{SourceQuotation}

ويرى الحاج، مع القديس أوغسطين، أن الشر المطلق مستحيل، لأنه لا يزال يرتقي إلى الخير. ويعد نظرية الإله المحسن أو المسيء خيالًا عاطفيًا محضًا، يناقضه العقل الإنساني ومظهر العالم. فكثيرًا ما يكون الشر الصورة الفاعلة للخير؛ وكما يقول ف و نيومان: «كذلك يكون الشر كشفًا عن الخير». والكائنات عنده كلها متساوية؛ فما دامت تملك «أغاسا» الهندوسية، سائل الحياة أو القوة الحيوية، فلا يهم أكانت:

\begin{SourceQuotation}
فطرًا أم سنديانة أم دودة أم إنسانًا.
\end{SourceQuotation}

ويقول إن الحرب تجلب مصائب فردية لا تُحصى، لكنها تدفع التقدم العام برفع الأقوى على أنقاض الأجناس الأضعف. وتخرّب الزلازل والأعاصير مناطق صغيرة؛ لكن الأولى تبني الأرض ليقيم عليها الإنسان، والثانية تجعل الهواء صالحًا لأن يتنفسه. ومن هنا يردّد:

\begin{SourceQuotation}
العلة الكلية\\
لا تعمل بنواميس جزئية، بل بنواميس عامة.
\end{SourceQuotation}

وتلحق بنظرة رجل الكنيسة غير الأخلاقية إلى «الخطيئة الأصلية» النظرية غير العلمية القائلة إن الشر دخل العالم مع آدم ونسله. فلنسأل: ماذا كانت حال كرتنا في الأيام السابقة لآدم، حين كان طغاة الأرض، الزواحف الهائلة وغيرها من الوحوش، يعيشون في صراع دائم، وفي تدمير لا نملك الآن إلا أضعف أمثلته؟ وما الحال الفعلية لعالم المياه، حيث لا غاية للحياة إلا الموت، وحيث ناموس القتل هو ناموس التطور؟

سيتهم بعضهم الحاج بقلة التوقير، ويعدّونه «نائبًا للشيطان يجلس في كرسي الوباء». لكنه لا يقصد قلة التوقير. ومثل رجال من مقام أعلى بكثير، ينكرون الإلهي إنكارًا إلهيًا، يقول ما يفكر فيه الآخرون ويخفونه. وهو يرى، مع مؤلف «الدين فوق الطبيعي»، أننا «نكسب أكثر بما لا يقاس مما نخسر حين نتخلى عن الاعتقاد بحقيقة الوحي»؛ ويتطلع إلى يوم «يكون فيه الطغيان القديم قد تحطم، وتكون فوضى الانتقال قد انقضت». لكنه شرقي. وحين يكرر قول اليوناني «تذكّر ألا تؤمن»، فإنه يعني: اجتهد لتتعلم وتعرف؛ فالأفكار الصحيحة تقود إلى الأفعال الصحيحة. ومن الأبيات المزدوجة التي لم تُترجم في هذه القصيدة:

\begin{SourceQuotation}
من جميع طرق الحياة الآمنة، تظل الطريق الأكثر أمانًا هي الشك؛\\
بالإيمان ينال الرجال العالم المقبل، ومن غيره ينالون العالم الحاضر.
\end{SourceQuotation}

وهذا هو قول الإسباني:

\begin{SourceQuotation}
من أكثر الأشياء يقينًا، أن الشك أكثر يقينًا.
\end{SourceQuotation}

وهو شعور حديث نموذجي لعصر العلم النحاسي، الذي أعقب عصر العاطفة الذهبي. لكن الحاج يتابع:

\begin{SourceQuotation}
يقول الحكماء، وأنا أقول لك: كلا! تلقَّ جميع العقائد بإيمان متساوٍ؛\\
لا يزيد ولا ينقص؛ فالشك موت، وأكثر الناس إيمانًا أكثرهم حياة.
\end{SourceQuotation}

وهنا أيضًا دقة شرقية؛ فمن يؤمن بكل شيء على نحو متساوٍ وعام، يمكن أن يقال إنه لا يؤمن بشيء. وليست هذه النظرة الأوروبية البسيطة التي تجعل الشك الصادق يعدل اثنتي عشرة عقيدة. وهي تعارض مباشرة ذلك الكاتب المشهور الذي يرى أن الرجل البسيط الإيمان يساوي تسعة وتسعين ممن لا يتمسكون إلا بالمصالح الأنانية لفرديتهم الخاصة. ويعني هذا القول المعتم، إن كان يعني شيئًا، أن قوى الإنسان المسماة أخلاقية، أي الخيال وتصور المثل، يجب أن تتسلط على قوى الإدراك والتأمل؛ وهذا محض عبث! وقد أنتج توريكريماتا، المعروف أيضًا بتوركيمادا، الذي كان يفيض دموعًا صادقة وهو يأمر بإحراق ضحاياه أحياء؛ وأنتشييتا، صانع المعجزات في البرازيل، الذي قطع رأس هرطوقي اهتدى، لئلا يفقد روحه الخالدة إذا انتكس من حال النعمة.

لكن هذا المسلك من التأمل، الذي يَسِمه المتعصبون بـ«الشك والإنكار والهدم»، وهذا الارتياب الديني الجاد، وهذا السؤال الفضولي: «هل للمأثور العام أساس من الواقع؟»، وهذا التوق إلى أسرار المستقبل وخفاياه، وإلى غير المرئي والمجهول، مشتركة بين جميع الأجناس وفي كل عصر. حتى عند الرومان، الذين كان رجلهم المثال في أيام أغسطس هو هوراس، الفيلسوف والأبيقوري، نجد بروبرتيوس يسأل:

\begin{SourceQuotation}
أحكاية مختلقة تلك التي نزلت على الشعوب البائسة؛\\
ولا يمكن للخوف أن يتجاوز محرقة الجنازة؟
\end{SourceQuotation}

لنعد إلى الموضوع: ستفاجئ مذاهب الحاج في الضمير والتوبة من لا يتابعون مسار فكره:

\begin{SourceQuotation}
لا تندم أبدًا لأن إرادتك لا تتحد مع إرادة القدر؛\\
فكّر، إن شئت، قبل أن تفعل، ولكن لا تأسَ أبدًا على الفعل بعد وقوعه.
\end{SourceQuotation}

وهذا أيضًا مذهبه المعدّل في القدر. وهو لا يقبل الطريقة الصاخبة لقطع العقدة الغوردية، التي يقترحها البريطاني النبيل ضيق الأفق: «نعرف أننا أحرار، وانتهى الأمر!» بل يؤثر قول لامارك: «الإرادة، في الحقيقة، ليست حرة أبدًا». ويعتقد أن الإنسان عنصر متسق في التقدم العظيم للطبيعة؛ ونتيجة لتفاعل الكائن وبيئته، يعمل خلال فترات من الزمان الكوني. وينظر إلى الآلة الإنسانية، أنبوب اللحم، على أنها تقوم على النظرية الفيزيائية للحياة. فكل واقعة وظاهرة جسدية تنمو، كالشجرة، من الداخل أو الخارج، ليست إلا نتاجًا للتنظيم؛ إذ تخضع الأجسام الحية للناموس الطبيعي الذي يحكم غير الحي وغير العضوي. وبينما يؤكد المتديّن أن الإنسان ليس مجرد لعبة للقدر، بل فاعل حر مسؤول أمام نفسه، وله عمل ينجزه وواجبات يؤديها، يرى الحاج، مع مدارس حديثة كثيرة، أن الذهن كلمة تصف عملًا مخصوصًا للمادة؛ وأن القوى عمومًا مظاهر لحركات في الجهاز العصبي المركزي؛ وأن كل فكرة، حتى فكرة الألوهية، ليست إلا نبضة صغيرة معينة في كتلة صغيرة معينة من عجينة حيوانية، هي الدماغ. ومن ثم لا يعترض على قرابته لقرد شبيه بالإنسان، نازل المنخرين، عديم الذيل، متحدّر من كائن أحادي الخلية أو من زقيّ بدائي.

ولهذا يقول في الواقع: «أتيت إلى العالم من غير أن أطلب إذنًا أو أحصل عليه؛ بل من غير أن يُطلب إذني أو أعطيه. وها أنا أجد نفسي مكتوف اليدين بالشروط، مقيّدًا بالقوانين والظروف التي لم يكن لكلمتي أي دور في صنعها. كنت في الرحم جهازًا يتحرك من تلقاء نفسه؛ وسيجدني الموت مجرد آلة. ولذلك فلست أنا، بل الناموس، أو واضع الناموس إن شئتم، مسؤولًا عن جميع أفعالي». ولألاحظ هنا أن «الناموس» يفترض في العقل الغربي واضعًا للناموس؛ وليس الأمر كذلك عند الشرقي، ولا سيما الصوفي الذي يعد هذه الأفكار بشرية، وقد وُسّعت بغير حق لتفسير غير البشري، الذي يسمّيه الناس الإلهي.

ويقول أيضًا: «أنا فرد، أي ما ليس فيه جزء قابل للقسمة، \textenglish{qui nil habet dividui}؛ دائرة تمس جيرانها وتتقاطع معهم عند نقاط معينة، لكنها لا تتطابق معهم في أي موضع ولا تمتزج بهم. جسديًا لست مطابقًا للآخرين في جميع النقاط. وأخلاقيًا أختلف عنهم؛ ولا تشبه مداخل المعرفة عندي، أي أعضاء حواسي الخمسة مع تأويلها على طريقة شيلّي، مداخل أي كائن آخر شبهًا تامًا. ومن ثم لن يكون أثر العالم والحياة والأشياء الطبيعية في حالتي هو نفسه عند أكثر الكائنات شبهًا بي. ولذلك أطالب بحق إنشاء النظام الذي يهمني أكثر من غيره، أو تعديله، لاستخدامي الخاص؛ وإذا رُفض هذا الإذن المعقول، أخذته بغير إذن.

«لكن فرديتي، مهما كانت كافية كل الكفاية لي، نقطة متناهية الصغر، وذرة تخضع في كل شيء لناموس العواصف الذي يدعى الحياة. أشعر وأعرف أن القدر موجود. لكن لا أستطيع أن أعرف ما قُدّر أن يصيبني وما لم يُقدّر. ولذلك تكمن في طلب الكمال بوصفي فردًا أسمى واجباتي، بل واجبي الوحيد، مع مزج أنا بنحن على الوجه الواجب. وأعترض على أن أكون إنسانًا بلا ذات؛ فهذا يدل عندي على حس أخلاقي مقلوب. وأنا ملزم بالتفكير بعناية في نتائج كل كلمة وفعل. لكن حين يصير المستقبل ماضيًا، سيكون محض عبث أن أحزن أو أتوب مما قضى به الناموس الكلي».

والاعتراض المعتاد هو ما يتعلق بممارسة الإنسان. فيقول: «هذا حسن نظريًا؛ لكن كيف يُنفّذ؟ لماذا، مثلًا، تقتل الرجل الذي أجبره القدر على قتل أبيك، أو تسلّمه ليُقتل؟» ويجيب الحاج عبدو: «أفعل كما يفعل الآخرون، لا لأن القتل وقع منه، بل لأنه لا ينبغي أن تُتاح للقاتل فرصة أخرى للقتل. إنه نمر ذاق الدم، وينبغي إطلاق النار عليه. وأنا مقتنع أنه كان أداة في يد القدر؛ لكن هذا لن يمنعني من اتخاذ التدابير، أكانت مقدّرة من قبل أم لم تكن، كيلا يُستعمل على النحو نفسه مرة أخرى».

وكما يكون الأمر مع التوبة يكون مع الضمير. فقد يكون الضمير «خوفًا هو ظل العدالة»، كما أن الشفقة ظل الحب. ومع أنه ليس إلا مصادفة جغرافية وزمنية تتبدل مع كل عصر من عصور العالم، فقد يردع الناس عن طلب ثمرة النذالة الناجحة ونيلها. لكن هذا الحافز إلى الإحسان ينبغي أن يُطبّق على الأفعال التي ستقع، لا على الأفعال التي وقعت.

ويميز الحاج أيضًا بعناية بين عمل القدر تحت حكم إله شخصي، وعمله تحت سيادة الناموس. ففي الحالة الأولى يكون التناقض بين العلم السابق للخالق وحرية إرادة المخلوق مباشرًا ومحسوسًا ومطلقًا. ويمكننا بالقدر نفسه أن نتحدث عن بياض أسود وسواد أبيض. لم تستطع مئة جيل من رجال اللاهوت حل اللغز؛ وسيفشل مليون جيل. فالصعوبة لا تُذلّل عند المؤمن بإله شخصي، الذي لا بد أن يكون إلهه القادر على كل شيء محيط العلم، ومن حيث هو محيط العلم عالِمًا مسبقًا. لكنها تزول حين نحوّل الشخص إلى ناموس، أو إلى ترتيب مستقر للأحداث، خاضع أيضًا لاستثناءات معينة ثابتة لا تتبدل، ولكن الإنسان لا يعرفها في الوقت الحاضر. والفرق جوهري، كالفرق بين القانون الجزائي بحظره الضيق، وبين الوصية الواسعة التي تهدي أكثر مما تتسلط.

وهكذا أيضًا يعدّل الاعتقاد بناموس ثابت، مقابل إرادة اعتباطية، آراء الحاج في طلب السعادة. فالبشر، «حيوان العلل الذي لا يهدأ»، \textenglish{das rastlose Ursachenthier}، يولدون ليكونوا في مجموع الأمر متساوين في السعادة والشقاء. وأعلى الكائنات تنظيمًا، خزف عائلتنا الرفيع، تتمتع أكثر وتتألم أكثر؛ فلها قدرة على الارتفاع إلى أعلى سماء اللذة، وعلى الغوص عميقًا في نهر الشقاء والألم السريع الجريان. وهكذا يقول دانتي في الجحيم، النشيد السادس، البيت ١٠٦:

\begin{SourceQuotation}
علمك الذي يريد\\
أنه كلما كان الشيء أكثر كمالًا؛\\
زاد إحساسه بالخير، وكذلك بالألم.
\end{SourceQuotation}

وهكذا تقرر البوذية أن الوجود في ذاته يستلزم الجهد والألم والحزن؛ وكلما ارتفع الكائن زاد ألمه. أما الطين العادي فيتمتع قليلًا ويتألم قليلًا. اجمع الكل ووزّع المجموع؛ تحصل على متوسط، ويكون المتسول، في مجموع الأمر، سعيدًا كالأمير. فلماذا إذن، يسأل المعترض، لا يكف الإنسان عن الكد والصراع كي يتغير ويرتفع، وهو صراع يتضمن فكرة تحسين حاله؟ ويجيب الحاج: «لأن ذلك هو الناموس الذي يولد الإنسان في حكمه؛ وقد يكون شرسًا كالمجاعة، وقاسيًا كالقبر، لكن على الإنسان أن يطيعه طاعة عمياء». ولا يدخل في مسألة ما إذا كانت الحياة تستحق أن تُعاش، أو ما إذا كان ينبغي للإنسان أن يختار أن يولد. ومع ذلك فإن تشاؤمه الشرقي، الذي يقابل بحدة تفاؤل الغرب، يردّد هذه الأبيات:

\begin{SourceQuotation}
حياة شحيحة بالنتائج الكبيرة؛\\
وإن أمكن احتمالها، لا تكاد تبدو جديرة\\
بكل هذا البذخ في الكلمات؛\\
ولا بألم الميلاد هذا.
\end{SourceQuotation}

الحياة، أيًا تكن نتيجتها، مبنية على أساس من الحزن. فالأدب، صوت الإنسانية، وحكم البشر، يعلنان أن كل وجود حال من الأسى. ويحاول «أطباء الروح» إنقاذ كآبتها من الانحدار إلى اليأس، بجرعات من الاعتقاد الثابت بحضور الله، وبضمان الخلود، وبرؤى الظفر الأخير للخير. ولو كان الحاج عبدو مجرد لاهوتي، لأضاف أن الخطيئة، لا إمكان العصيان بل العصيان نفسه على الضمير، هي الصورة الأولى للشر، لأنها تنتج الخطأ الأخلاقي والعقلي. فهذا الإنسان الذي يهمل قراءة ناموس الضمير، مهما اختلف عن قانون المجتمع، مذنب بالإهمال. وذاك الذي يعتم نور الطبيعة بالسفسطة، يصير عاجزًا عن تمييز حقائقه الخاصة. وفي الحالتين يعقب الخطأَ المتبنّى عمدًا ألمٌ يُقال لنا إنه يأتي عدلًا وإحسانًا، إنذارًا وعلاجًا وتأديبًا.

لكن الحاج لا يرضى عن فكرة أن الشر ينشأ من الأفعال الفردية لفاعلين أحرار، نحن والآخرين. فهذا المذهب لا يفسر خصائصه: جوهريته وعموميته. وقد يبدو طبيعيًا ألا تختار دائمًا الخيرَ كائناتٌ لا تملك إلا إمكان الحرية. ولكن ألا يوجد بين مليارات البشر الذين سكنوا الأرض واحد يختار الخير بلا استثناء، يبرهن على عدم كفاية هذا الحل. ولهذا لا يؤمن أحد بوجود الإنسان الكامل في الحال الراهنة للأشياء. ويرفض الحاج جميع التفسيرات الشائعة والأسطورية بسقوط «آدم»، والفساد الفطري للطبيعة الإنسانية، والكمال المطلق لبعض التجسدات الذي يُستدل به على ألوهيتها. ولا يسعه إلا أن ينوح على شيوع الشر، ويفترض أن أساسه الخطأ، ويقصد إلى تخفيفه باجتثاث الجهل الذي يلده ويغذيه.

ومبحثه في المآل الأخير، مثل مبحث الصوفية عمومًا، غامض شبحي. وقد يميل إلى مذهب ماركوس أوريليوس: «العنب غير الناضج والناضج واليابس؛ كل الأشياء تغيرات، لا إلى العدم، بل إلى ما ليس موجودًا في الحاضر». وهذه واحدة من «الطفرات الوحشية المروّعة للآراء»، \textenglish{monstruosa opinionum portenta}، التي ذكرها المجمع العام التاسع عشر، المسمى أيضًا مجمع الفاتيكان الأول. لكنه لا يقبلها إلا بقيد. فهو يتمسك بالعبادة الأخلاقية للطبيعة، لا بعبادتها العقلية؛ والطبيعة يعرّفها المحدثون بأنها «مرادف غير علمي ومتخيَّل للمجموع الكلي للظواهر المرصودة». ولهذا يتمسك بـ«المذاهب المظلمة المحطّة للمادي»، المؤلّه للمادة، في مقابل الروحاني؛ وهو تمييز أشد بروزًا بكثير في الغرب منه في الشرق. فأوروبا ترسم خطًا قاسيًا جافًا بين الروح والمادة؛ وآسيا لا تفعل.

ويعترض المثالي عندنا على الماديين بأنهم لا يستطيعون الاتفاق على النقاط الأساسية؛ ولا تعريف الذرة؛ ولا تفسير تحول الفعل الفيزيائي والحركة الجزيئية إلى وعي، أو العكس؛ ولا القول ما المادة؛ وأخيرًا بأن بيركلي ومدرسته أثبتوا وجود الروح مع إنكار وجود المادة.

ويجيب الماديون بأن غياب الاتفاق لا يدل إلا على دراسة لم تتقدم بما يكفي؛ وبأن الإنسان لا يستطيع وصف الذرة لأنه لا يزال طفلًا في العلم، ولكن لا سبب يمنع رجولته الناضجة من عبور الخطأ والعجز إلى الحقيقة والمعرفة؛ وبأن الوعي يصير خاصية للمادة حين تتوافر شروط معينة؛ وبأن الهيولى، \textenglish{Hyle}، أو المادة، يمكن أن تُعرّف مؤقتًا بأنها «ظواهر لها بنية تحتية تخصها، متعالية وأبدية، خاضعة للفعل المباشر أو غير المباشر للحواس الخمس، فيما تتمثل خصائصها في ثلاث حالات: الصلبة والسائلة والغازية». وهم يفضلون الحس السليم للبشر على بيركلي كثير المحاجّة. ويسألون المثالي والروحاني لماذا لا يستطيعان العثور على اسمين لنفسيهما من غير الاقتراض من مدرسة «مظلمة منحطّة»؛ ولماذا يجب أن يسمّي الأول نفسه نسبة إلى عينه، \textenglish{idein}، والثاني نسبة إلى نفَسه، \textenglish{spiritus}. ومن ثم يعيّرهم الحاج بإلصاق حدودهم الخاصة بقدرتهم الكلية، وبإنزال السماء إلى السوق، كما قال سقراط.

ويميل الفكر الحديث أكثر فأكثر إلى رفض المثالية الخام، وإلى تأييد النظرية الواحدية، ونظرية الوجهين، والواقعية المتحوّلة. وهو يناقش طبيعة الأشياء في ذاتها. وعلى سؤال: هل يوجد خارجنا شيء يقابل إحساساتنا، أي هل العالم كله مجرد «أنا»؟ يجيبون بأن شيئًا آخر موجود بوضوح؛ وأن هذا الشيء الآخر يحدث اضطراب الدماغ الذي يسمّى الإحساس. وتأمرنا الغريزة أن نفعل شيئًا؛ ويوجّه العقل، وهو ميزان القوى، ويسيطر الدافع الأقوى. ويبدو أن العلم الحديث، باكتشاف «المادة المتألقة»، وهي حالة رابعة، يوفّق بين المدرستين. ويقول أحد أصحاب المراجعات: «إن اكتشاف حالة رابعة للمادة هو باب مفتوح على لا تناهي تحوّلاتها؛ وهو إمكان أن يكون الإنسان غير مرئي وغير ملموس مع بقائه ذا جوهر؛ وهو دخول عالم الأرواح مجال الفرضيات العلمية من غير عبث؛ وهو إمكان أن يؤمن المادي بحياة وراء القبر، من غير أن يتخلى عن الأساس المادي الذي يراه لازمًا لحفظ الفردية».

الروح عند الحاج عبدو ليست مادية؛ إذ سيكون ذلك تناقضًا في الألفاظ. وهو يعدّها، مع محدثين كثيرين، حالًا من الأحوال، لا شيئًا؛ وكلمة ملائمة تدل على الإحساس بالشخصية وبالهوية الفردية. وفي معناها الشبحي يجد عقيدة مصطنعة، يكاد يستحيل أن تكون للهمج الوحشيين في العصر الحجري. ويجدها في الكتب الجنائزية لمصر القديمة، التي ربما انتقلت منها إلى الزند أفستا والفيدات. وهي غير معروفة في أسفار التوراة العبرية الخمسة، التي ما زال يُنسب قسم منها إلى موسى؛ أو، بالأحرى، يتجاهلها المؤلف أو المؤلفون عمدًا. ولم يستطع المسيحيون الأوائل الاتفاق في الموضوع؛ فقد دعا أوريجانوس إلى وجود أرواح البشر السابق، مفترضًا أنها خُلقت جميعًا في وقت واحد، ثم تجسدت بالتتابع. ويجعل آخرون الروح تولد في ساعة الميلاد؛ وهكذا.

لكن فعل الدماغ، أو الذهن إن شئتم تسميته كذلك، لا يقتصر على قوى الاستدلال؛ ولا يسعنا تجاهل المشاعر والعواطف، التي لعلها أقوى حقائق الحياة. وصوتها المثبت العالي يقابل بقوة نبرات الفكر المتلعثمة. ويبدو أنها تطلب حياة مقبلة، بل حالًا من الثواب والعقاب من صانع العالم: البستاني الأبدي، \textenglish{Ortolano Eterno}، وفخّار الشرق، وصانع ساعات الغرب. وهي تحتج على فكرة الفناء، وتثور على تصور الفراق الأبدي للوالدين والأقارب والأصدقاء. ومع ذلك ليست عقيدة الحياة المقبلة عامة جامعة بأي حال. ويبدو أن الجنس الإنجليزي الأوروبي لا يستطيع الوجود من غيرها؛ وقد سمعنا مؤخرًا عن «أرض الروح الآرية». ومن جهة أخرى تبشّر مدارس بوذية، وحتى برهمانية، كثيرة بالنيرفانا، أي عدم وجود نسبي، وبالبارينيرفانا، أي عدم مطلق. كذلك فإن العائلة الطورانية الكبرى، التي تشغل بالفعل جميع شرقي آسيا، ظلت تتجاهلها؛ ويحتج مئتا مليون كونفوشيوسي صيني، وهم معظم الأمة، بقوة على الدعامة الرئيسية للعقائد الغربية، لأنها «تجعل الناس غير صالحين لعمل الحياة وواجبها، بتثبيت تأملاتهم في عالم مجهول». وحتى أتباعها، في كل العصور والأجناس والعقائد، لا يستطيعون إنكار أن العالم المقبل نسخة، أُضفي عليها من المثالية قدر يزيد أو ينقص، من العالم الحاضر؛ وأنه يفتقر إلى تفصيل واحد تشتم منه رائحة الأصالة. إنه في الواقع مجرد استمرار؛ وهذا الاستمرار «لم يثبت».

\begin{SourceQuotation}
ما أصعب أن يكون المرء إنسانًا!
\end{SourceQuotation}

وعزاء الحاج الوحيد هو تنمية النفس ومتع العواطف. وقد يكون هذا التعاطف حبًا غير مباشر للذات، وانعكاسًا لنور الأنانية؛ لكنه انتقل على نحو يستلزم مذهبًا آخر من القناعات. وهو يحتاج إلى اسم مختلف؛ فتسمية الإحسان «حبًا للذات» تعني أن الثمرة أو الزهرة لا تعتمد على جذر لنموها فحسب، وهذا صحيح، بل إنها الجذر نفسه، وهذا خطأ. وأخيرًا فإن مثاله من أعلى المثل؛ وهو يدّخر مديحه لـ:

\begin{SourceQuotation}
حيوات عِيشت في طاعة الناموس الباطن؛\\
الذي لا يستطيع أن يتبدل.
\end{SourceQuotation}

البستاني الأبدي؛ كما يقول النقش القديم:

\begin{SourceQuotation}
الإنسان وُضع في بستان؛\\
وحُكم عليه في بستان؛\\
ودُفن في بستان؛\\
ووُلد من جديد في بستان.
\end{SourceQuotation}
